\chapter{Vita media di un protone}

In base alle conoscenze attuali di fisica delle particelle il protone libero è una particella stabile. Ciò vuol dire che esso non decade in altre particelle.\\
Questo è dovuto alla conservazione del numero barionico\footnote{Il numero barionico è un numero quantico additivo che si conserva in tutti i processi finora osservati.} nei processi elementari. Infatti il barione\footnote{I barioni sono le particelle composte da tre quark. Il protone è il barione più leggero: se decadesse conservando $B$ dovrebbe produrre almeno un barione nello stato finale, che però sarebbe più pesante, per cui il $Q$-valore risulterebbe negativo e il decadimento energeticamente vietato.} più leggero è proprio il protone.\\
Tuttavia, alcuni modelli teorici di grande unificazione (GUT), come il modello di Georgi--Glashow, prevedono processi di non conservazione del numero barionico, tra cui il decadimento del protone. Analizzando il decadimento protonico sarebbe quindi possibile ricavare informazioni sulle teorie di grande unificazione.

Per questo motivo nel mondo sono attivi diversi esperimenti che hanno come obiettivo quello di misurare la vita media del protone. Tale evento, però, se esiste è estremamente difficile da osservare e richiede apparati molto grandi e complessi.\\

Di seguito si riportano alcuni ragionamenti che permettono di porre un limite inferiore alla vita media ($\tau_p$) del protone. 

\section{Attività del corpo umano}

Il primo argomento non richiede alcun apparato sperimentale ed è basato su una sola osservazione: noi siamo vivi. I protoni del nostro corpo, infatti, se decadessero rilascerebbero energia nel corpo stesso, e se lo facessero troppo rapidamente la dose accumulata nel corso di una vita sarebbe letale. Il fatto che una persona possa vivere cento anni pone quindi un limite inferiore a $\tau_p$.\\

Nel corpo umano troviamo atomi, come $^{40}\mathrm{K}$ e $^{14}\mathrm{C}$, i cui nuclei sono soggetti a decadimenti: è possibile stimare che in media l'attività del corpo umano sia dell'ordine di $A = 5 \cdot 10^3 \div 10^4$ Bq. A questa radioattività naturale si aggiungerebbe, se il protone fosse instabile, quella dovuta al suo decadimento. 

Quando l'energia così assorbita diventa pericolosa per il corpo? Dipende dalla dose\mkeyword{dose} ($D$), definita come
%
\begin{equation}
    D = \frac{E_{ass}}{M} \quad , \quad [D] = \text{Gy} = 1 \ \text{J}/\text{kg}
\end{equation}
dove $E_{ass}$ è l'energia assorbita dal corpo mentre $M$ è la sua massa. Una dose di circa $1$ Gy a tutto il corpo corrisponde alla soglia letale.\\

Supponendo che un uomo viva circa $100$ anni e che $D_{max}$ sia la dose massima che il corpo può assorbire, si dovrà avere 
%
\begin{equation}
    D_{max} > D(100 \ y)
    \label{eq:disuguaglianza_dose}
\end{equation}
Prendendo la dose come una funzione del tempo e considerando solo i decadimenti dei protoni, si ha 
%
\begin{equation}
    D(t) = \frac{E_{dec} \cdot N_{dec}(t)}{M}
\end{equation}
dove $E_{dec}$ è l'energia rilasciata da un protone in un decadimento, la quale corrisponderà a una frazione dell'energia totale del protone
%
\begin{equation*}
    E_{dec} = f \cdot m_p c^2 \quad , \quad f \in (0,1]
\end{equation*}
Mentre $N_{dec}$ è il numero medio di protoni decaduti all'istante $t$, che sappiamo essere 
%
\begin{equation*}
    N_{dec}(t) = N_p \left(1 - e^{-t/\tau_p}\right)
\end{equation*}
dove $N_p$ è il numero di protoni nel corpo. Indicando con $\gamma_p$ la frazione di nucleoni che sono protoni, che vale $0,5 \div 0,55$, e trascurando la massa degli elettroni, il numero di nucleoni è $N_{\text{nucleoni}} = M/m_u$, da cui
%
\begin{equation*}
    N_p = \gamma_p \, N_{\text{nucleoni}} = \gamma_p \frac{M}{m_u} \quad , \quad m_u = 1,66 \cdot 10^{-27} \ \text{kg}
\end{equation*}
Inoltre, sapendo che la vita media di un protone è molto maggiore di $100$ anni, possiamo approssimare $1 - e^{-t/\tau_p} \approx t/\tau_p$, cosa che porta a scrivere il numero medio di protoni decaduti come 
%
\begin{equation*}
    N_{dec}(t) = \gamma_p \frac{M}{m_u}\cdot \frac{t}{\tau_p}
\end{equation*}

Sostituendo tutto nella funzione dose si ha 
%
\begin{equation*}
    D(t,\tau_p, f) = \frac{1}{M} \cdot \gamma_p \frac{M}{m_u} \cdot \frac{t}{\tau_p} \cdot f m_p c^2 = \frac{\gamma_p \cdot f m_p c^2}{ m_u} \cdot \frac{t}{\tau_p} = 5 \cdot 10^{16} \cdot f \cdot \frac{t}{\tau_p} \ \text{Gy}
\end{equation*}
Si noti che la massa del corpo si semplifica: il risultato \emph{non dipende} dalla massa della persona considerata, ma solo dalla composizione in nucleoni della materia di cui è fatta. Il limite che otterremo vale quindi per chiunque.\\

Andiamo adesso a sfruttare la disuguaglianza \eqref{eq:disuguaglianza_dose}. Ponendo $D_{max} = 1$ Gy si ottiene
%
\begin{equation*}
    D_{max} > D(100 \ y, \tau_p, f) \implies \tau_p > 5f \cdot10^{18} \ y
\end{equation*}

Resta da stimare $f$. Nel canale di decadimento dominante previsto dalle GUT, $p \to e^+ + \pi^0$, il positrone annichila con un elettrone del mezzo e il pione decade in due fotoni: tutti i prodotti depositano quindi la loro energia nel corpo, e praticamente tutta l'energia di massa del protone viene assorbita. Solo i canali che producono neutrini sottraggono energia al bilancio. Possiamo quindi assumere $f \sim 1$ e concludere, 

\begin{tcolorbox}[colback=black!5!white, colframe=black!45!white, boxrule=0.8pt]
\begin{equation}
    \tau_p > 10^{18} \ y
\end{equation}
\end{tcolorbox}

È notevole che un limite di questa forza si ottenga senza alcun apparato sperimentale, semplicemente osservando che siamo vivi, tuttavia esistono esperimenti che permettono di andare a porre un limite molto più forte sul valore di $\tau_p$.

\section{Conteggio diretto dei decadimenti}

Andiamo adesso ad allestire un vero e proprio esperimento per porre un limite ancora più stringente sul valore della vita media del protone. \\

Supponiamo di avere $N_p = 8 \cdot 10^{32}$ protoni e di osservare i decadimenti di questi per un tempo pari a $T = 5 \ y$ e che la sensibilità dello strumento, cioè la probabilità di vedere quando un protone decade, sia $\varepsilon = 40 \%$. Supponiamo inoltre che \emph{non ci sia fondo}\mkeyword{fondo}. Si dice fondo l'insieme degli eventi che il rivelatore registra e che assomigliano a un decadimento del protone, pur essendo prodotti da altri processi: le interazioni dei neutrini atmosferici, i muoni cosmici e la radioattività naturale dei materiali di cui è fatto l'apparato. Dire che il fondo è nullo significa che nel tempo di osservazione ci si aspetta zero eventi di questo tipo, e quindi che qualunque conteggio osservato sarebbe necessariamente un decadimento di un protone\footnote{È un'idealizzazione: negli esperimenti reali il fondo è piccolo ma non nullo, ed è dominato dai neutrini atmosferici, che attraversano qualunque schermatura. Gli altri contributi si riducono collocando il rivelatore in profondità sotto roccia e selezionando materiali a bassa radioattività.}.\\

A seguito dell'esperimento si hanno \emph{zero} conteggi. 

Ricordando che il decadimento è un processo poissoniano, la probabilità di vedere $k$ conteggi in un tempo $T$ risulta essere 
%
\begin{equation*}
    \mathcal{P}(k|\mu) = \frac{\mu^k e^{-\mu}}{k!}
\end{equation*}
Ponendo $k = 0$ si ottiene 
%
\begin{equation*}
    \mathcal{P}(0|\mu) = e^{-\mu}
\end{equation*}
Dove 
%
\begin{align*}
    \mu_{teorico} = N_p(1 - e^{-T/\tau}) \approx N_p \frac{T}{\tau_p}\\
    \implies \mu = \varepsilon \mu_{teorico} = N_p \varepsilon \frac{T}{\tau_p}
\end{align*}
La quantità $\mu$ è quindi il numero medio di conteggi che ci aspetteremmo di osservare, e dipende dal valore di $\tau_p$, che è l'incognita del problema. L'esperimento ha invece fornito un risultato preciso: zero conteggi. 

Andiamo a vedere quali valori di $\tau_p$ sono compatibili con questa osservazione. 

Fissiamo una soglia di probabilità e dichiariamo esclusi tutti i valori di $\tau_p$ che renderebbero la nostra osservazione più improbabile di quella soglia. Scegliendo il $90\%$, convenzionale in fisica delle particelle, escludiamo tutti i valori per cui la probabilità di non osservare nulla sarebbe inferiore al $10\%$. 

Chiamando $\mu_{90}$ il valore limite si ha quindi
%
\begin{equation*}
    \mathcal{P}(0|\mu_{90}) = e^{-\mu_{90}} = 0,1 \qquad \Longrightarrow \qquad \mu_{90} = \ln 10 \simeq 2,3
\end{equation*}
Equivalentemente: se il numero medio atteso fosse $\mu_{90}$, avremmo avuto il $90\%$ di probabilità di osservare \emph{almeno} un conteggio, dato che $\mathcal{P}_{oss} = 1 - e^{-\mu}$. Non avendone osservato nessuno, concludiamo
%
\begin{equation*}
    \mu < \mu_{90} \qquad \Longrightarrow \qquad N_p \, \varepsilon \, \frac{T}{\tau_p} < 2,3
\end{equation*}
e, siccome $\mu$ è decrescente in $\tau_p$, la disuguaglianza si rovescia dando un limite \emph{inferiore} sulla vita media

\begin{tcolorbox}[colback=black!5!white, colframe=black!45!white, boxrule=0.8pt]
\begin{equation}
    \tau_p > \frac{N_p \, \varepsilon \, T}{2,3} \qquad (90\% \ \text{CL})
    \label{eq:limite_poissoniano}
\end{equation}
\end{tcolorbox}
Sostituendo i valori del nostro esperimento si ottiene
%
\begin{equation*}
    \tau_p > \frac{8 \cdot 10^{32} \cdot 0,4 \cdot 5 \ y}{2,3} \simeq 7 \cdot 10^{32} \ y
\end{equation*}
cioè quattordici ordini di grandezza in più rispetto al limite ricavato dalla dose assorbita dal corpo umano.\\

\section{Super-Kamiokande}

L'esperimento che ad oggi fornisce i limiti più stringenti è Super-Kamiokande\mkeyword{Super-Kamiokande}, situato in Giappone a circa $1000$ m di profondità sotto il monte Ikenoyama. Si tratta di un cilindro di $39$ m di diametro e $41$ m di altezza riempito con $50$ kton di acqua purissima, osservata da oltre $11\,000$ fotomoltiplicatori montati sulle pareti. 

La scelta dell'acqua risponde a tre esigenze contemporaneamente: è economica, cosa necessaria viste le masse in gioco, contiene moltissimi protoni, ed è trasparente, il che permette di rivelare la luce Cherenkov emessa dalle particelle cariche prodotte nel decadimento. Ogni particella carica relativistica genera un cono di luce che i fotomoltiplicatori registrano come un anello: dalla forma e dall'intensità degli anelli si ricostruiscono il tipo, l'energia e la direzione delle particelle. 

I protoni disponibili sono di due tipi. Quelli dei nuclei di idrogeno sono \emph{liberi} e a riposo, mentre quelli legati nell'ossigeno si muovono con il moto di Fermi all'interno del nucleo e i prodotti del loro decadimento possono riurtare contro gli altri nucleoni prima di uscire. I primi danno quindi un segnale molto più pulito.\\

Il canale di decadimento principale cercato è
%
\begin{equation*}
    p \longrightarrow e^+ + \pi^0 \qquad \text{con} \qquad \pi^0 \to \gamma \gamma
\end{equation*}
che produce tre anelli Cherenkov, uno del positrone e due dei fotoni. La firma che distingue questo evento dal fondo è puramente cinematica, e discende dal fatto che il protone decade sostanzialmente a riposo: la massa invariante ricostruita di tutti i prodotti deve valere circa $m_p$, e l'impulso totale deve essere circa nullo, cioè le particelle devono uscire ``schiena contro schiena''. Richiedendo $800 < M_{tot} < 1050 \ \mathrm{MeV}/c^2$ e $P_{tot} < 250 \ \mathrm{MeV}/c$ si eliminano quasi tutte le interazioni dei neutrini atmosferici, che non rispettano questi vincoli.\\

I risultati dell'analisi su tutti i dati raccolti tra il $1996$ e il $2018$\footnote{Super-Kamiokande Collaboration, A. Takenaka \emph{et al.}, \emph{Search for proton decay via $p \to e^+\pi^0$ and $p \to \mu^+\pi^0$ with an enlarged fiducial volume in Super-Kamiokande I-IV}, Phys. Rev. D \textbf{102}, 112011 (2020).} sono riassunti nella tabella seguente. 

\begin{table}[htbp]
\centering
\begin{tabular}{lc}
\hline
Esposizione & $450$ kton$\cdot$anno \\
Efficienza di rivelazione & $\simeq 40 \%$ \\
Fondo atteso & $0,59$ eventi \\
Candidati osservati & $0$ \\
\hline
\end{tabular}
\caption{Risultati della ricerca del decadimento $p \to e^+ \pi^0$ in Super-Kamiokande.}
\label{tab:superk}
\end{table}

Si noti che il fondo atteso è inferiore all'unità: dopo vent'anni di presa dati l'ipotesi di fondo nullo fatta nella sezione precedente è ancora una buona approssimazione, ed è esattamente ciò che la profondità del rivelatore e la selezione cinematica servono a garantire. 

Possiamo allora applicare direttamente la \eqref{eq:limite_poissoniano}. Un kton d'acqua contiene circa $3,3 \cdot 10^{32}$ protoni, per cui l'esposizione corrisponde a $N_p T \simeq 1,5 \cdot 10^{35}$ protoni$\cdot$anno, da cui
%
\begin{equation*}
    \tau_p > \frac{N_p \, \varepsilon \, T}{2,3} \simeq \frac{1,5 \cdot 10^{35} \cdot 0,4}{2,3} \simeq 3 \cdot 10^{34} \ y
\end{equation*}
in ottimo accordo con il limite pubblicato dalla collaborazione

\begin{tcolorbox}[colback=black!5!white, colframe=black!45!white, boxrule=0.8pt]
\begin{equation}
    \frac{\tau_p}{B(p \to e^+\pi^0)} > 2,4 \cdot 10^{34} \ y \qquad (90\% \ \text{CL})
\end{equation}
\end{tcolorbox}

Due osservazioni per concludere. 

La prima riguarda la forma del risultato: quello che si misura non è la vita media totale, ma la vita media \emph{parziale} $\tau_p/B$ relativa a quel canale, dove $B$ è la frazione di diramazione. È la stessa distinzione incontrata parlando dei branching ratio nei decadimenti nucleari, ed è inevitabile: un esperimento può escludere solo i decadimenti che è in grado di riconoscere. 

La seconda è che un limite di questa entità ha già conseguenze teoriche. Il modello di Georgi--Glashow prevedeva $\tau_p \sim 10^{31}$ anni, tre ordini di grandezza al di sotto: la mancata osservazione del decadimento lo ha di fatto escluso, e le teorie di grande unificazione ancora in gioco sono quelle che prevedono vite medie più lunghe. Non aver visto nulla, in questo caso, è stata un'informazione preziosa.